\section*{Project contract}
Implement an ordinary and a tiled transpose for separate row-major word matrices.
Each dimension is between zero and sixteen inclusive. Source and destination active
windows fit in their separate 256-word allocations and may begin inside those
allocations. The pointers in \texttt{\$a0} and \texttt{\$a1} name those active
windows; \texttt{\$a2} and \texttt{\$a3} give the source row and column counts.
Return the copied count in \texttt{\$v0}, preserve all 64 bits of the saved registers
and restore \texttt{\$sp}. Do not perform I/O.

Copy each active source word once to its transposed destination, preserving its
32-bit pattern. Read that source cell before writing its destination. Matrix
accesses must produce aligned word notices. Leave input, guards, and inactive
output untouched, including writes later restored. Ordinary reads and writes each
follow source row-major order. The tiled function completes one clipped four-by-four
tile before accessing another. Tile order and register allocation may vary.
Auxiliary and stack traffic outside the matrix buffers is permitted under the
calling contract and contributes to whole-function cache counters.

Use the supplied stopped return boundary \texttt{transpose\_done}; its continuation
stores the count in \texttt{copied} and exits normally. The public checker invokes
both functions directly with varied arguments. Exact reference instruction and
cache counts illustrate a measurement; they are not submission requirements.
The downloadable assignment and public checks specify the complete executable
contract. Follow the guided edits or implement that contract in free-writing mode.
